\section*{Bab \ref{ch:b3:seismology-imaging} --- Seismologi dan Pencitraan Bumi}

\begin{solution}{exo:b3:seismology-imaging:1}
(a) $\sin i_0 = pv_0/R = 500\times5.80/6371 = 0.4552$, sehingga $i_0 = \ang{27.1}$.
(b) Pada \qty{771}{km}, $r = \qty{5600}{km}$ dan $r/v = 5600/11.066 =
\qty{506.1}{s.rad^{-1}}$; pada \qty{871}{km}, $r = \qty{5500}{km}$ dan $r/v =
5500/11.245 = \qty{489.1}{s.rad^{-1}}$. Sinar mencapai \omterm{def:b3:seismology-imaging:ray-parameter}{titik baliknya} saat $r/v = p = 500$:
$z_t = 771 + 100\times(506.1 - 500)/(506.1 - 489.1) = 771 + 36 =
\qty{807}{km}$.
\end{solution}

\begin{solution}{exo:b3:seismology-imaging:2}
(a) $(455.5 - 369.6)/10 = \qty{8.59}{s}$ per derajat, dan
$8.59\times180/\pi = \qty{492.2}{s.rad^{-1}}$. (b) Satu derajat busur adalah
$6371\pi/180 = \qty{111.19}{km}$, sehingga muka gelombang menyapu permukaan dengan laju
$111.19/8.59 = \qty{12.94}{km.s^{-1}}$, jauh lebih cepat daripada batuan mana pun di
permukaan: sinar tiba secara curam. (c) Di \omterm{def:b3:seismology-imaging:ray-parameter}{titik balik} $v = r_t/p =
5450/492.2 = \qty{11.07}{km.s^{-1}}$, yaitu kecepatan gelombang P di mantel bawah.
\end{solution}

\begin{solution}{exo:b3:seismology-imaging:3}
(a) $\mu_1 = \rho_1\beta_1^2 = 2700\times3500^2 = \qty{3.31e10}{Pa}$
(\qty{33.1}{GPa}); $\mu_2 = 3300\times4500^2 = \qty{6.68e10}{Pa}$
(\qty{66.8}{GPa}). (b) $\lambda = cT = 3.79\times20 = \qty{75.8}{km}$, 2.17
kali tebal lapisan: gelombang menjangkau jauh ke dalam setengah-ruang, dan hal inilah
yang menjelaskan mengapa kecepatannya sudah di atas $\beta_1$. (c) $4000/3.384 = \qty{1182}{s}$ dan
$4000/4.333 = \qty{923}{s}$: gelombang 20-s tiba \qty{259}{s}, atau
\qty{4.3}{min}, setelah gelombang 100-s.
\end{solution}

\begin{solution}{exo:b3:seismology-imaging:4}
(a) ${}_0S_2$: $1/(53.89\times60) = \qty{3.093e-4}{Hz} = \qty{0.3093}{mHz}$;
${}_0T_2$: $1/(44.01\times60) = \qty{0.3787}{mHz}$. (b) $1440/53.89 = 26.7$
dan $1440/44.01 = 32.7$ osilasi per hari. (c) Selisihnya adalah
$\qty{0.0694}{mHz}$, yang kebalikannya \qty{14400}{s}: setidaknya \qty{4.0}{h}
rekaman, dalam praktik jauh lebih panjang, karena puncaknya melebar akibat atenuasi dan
terbelah.
\end{solution}

\begin{solution}{exo:b3:seismology-imaging:5}
(a) Tali busur antara dua titik pada bola yang dipisahkan oleh sudut $\Delta$
di pusatnya memiliki panjang $2R\sin(\Delta/2)$, ditempuh dengan kecepatan $v$.
(b) Pada \ang{30}: $2\times6371\times\sin\ang{15}/10.0 = \qty{329.8}{s}$, dibandingkan
dengan \qty{369.6}{s}; pada \ang{90}: $2\times6371\times\sin\ang{45}/10.0 =
\qty{901.0}{s}$, dibandingkan dengan \qty{779.7}{s}. (c) Sinar menuju \ang{30} mencapai \omterm{def:b3:seismology-imaging:ray-parameter}{titik balik} di
mantel atas, tempat kecepatan sebenarnya sekitar 8 hingga \qty{10}{km.s^{-1}}, di bawah
\qty{10}{km.s^{-1}} yang seragam; sinar menuju \ang{90} menghabiskan sebagian besar lintasannya di
mantel bawah, dengan kecepatan 11 hingga \qty{13.7}{km.s^{-1}}, dan membelok ke bawah memasukinya, sehingga
Bumi yang sebenarnya lebih cepat.
\end{solution}

\begin{solution}{exo:b3:seismology-imaging:6}
(a) Dengan mengkuadratkan, $(2 - \xi)^4 = 16(1 - \xi/3)(1 - \xi)$. Kembangkan:
$16 - 32\xi + 24\xi^2 - 8\xi^3 + \xi^4 = 16 - \tfrac{64}{3}\xi + \tfrac{16}{3}\xi^2$,
sehingga $\xi^4 - 8\xi^3 + \tfrac{56}{3}\xi^2 - \tfrac{32}{3}\xi = 0$, yaitu
$\xi\,(3\xi^3 - 24\xi^2 + 56\xi - 32) = 0$, dan
$3\xi^3 - 24\xi^2 + 56\xi - 32 = (\xi - 4)(3\xi^2 - 12\xi + 8)$. Akar
$\xi = 0$ bukan gelombang. (b) Akar-akar persamaan kuadratnya adalah $2 \pm 2/\sqrt3$,
yaitu 3.155 dan 0.8453; hanya $\xi = 0.8453$ yang terletak di $(0, 1)$ (dan $\xi = 4$
tidak), dan ia memenuhi persamaan sebelum dikuadratkan, kedua ruasnya bernilai $4/3$.
Jadi $c/\beta = \sqrt{0.8453} = 0.9194$. (c) $c = 0.9194\times3.46 =
\qty{3.181}{km.s^{-1}}$; sepanjang \qty{1000}{km}, $314.4 - 289.0 = \qty{25.3}{s}$
setelah gelombang S.
\end{solution}

\begin{solution}{exo:b3:seismology-imaging:7}
$\mathsf G$ memiliki baris $(1, x_i)$ dan $\mathsf C_d^{-1} = \operatorname{diag}
(100, 100, 25, 25)$ (\unit{s^{-2}}). Persamaan normalnya adalah
\[
  \begin{pmatrix} 250 & 23750 \\ 23750 & \num{2812500} \end{pmatrix}
  \begin{pmatrix} a \\ b \end{pmatrix}
  = \begin{pmatrix} 3665 \\ 418125 \end{pmatrix},
\]
dengan determinan \num{1.390625e8}. Dengan demikian $b = (250\times418125 -
23750\times3665)/\num{1.390625e8} = \qty{0.12575}{s.km^{-1}}$ dan
$a = (2812500\times3665 - 23750\times418125)/\num{1.390625e8} =
\qty{2.713}{s}$. Kovariannya adalah $\mathsf N^{-1}$: $\sigma_b =
\sqrt{250/\num{1.390625e8}} = \qty{0.00134}{s.km^{-1}}$, $\sigma_a =
\sqrt{2812500/\num{1.390625e8}} = \qty{0.142}{s}$. Kecepatannya adalah
$1/b = \qty{7.95}{km.s^{-1}}$, dengan galat $\sigma_b/b^2 =
\qty{0.08}{km.s^{-1}}$: kecepatan di mantel paling atas. Residunya adalah $-0.001$,
$0.011$, $-0.076$ dan $\qty{0.036}{s}$, dan $\chi^2 = 0.19$ untuk $4 - 2 = 2$
derajat kebebasan: lebih kecil daripada yang diharapkan, sehingga galat yang dinyatakan
bahkan agak longgar; kecocokannya dapat diterima.
\end{solution}

\begin{solution}{exo:b3:seismology-imaging:8}
(a) $\mathsf G = (30\;\;40)$ berupa satu baris: $\mathsf G^{\mathsf T}\mathsf G$ adalah
berukuran $2\times2$ berpangkat satu, sehingga singular; model apa pun dengan $30m_1 + 40m_2 = 0$ dapat
ditambahkan tanpa mengubah datanya. (b) $\mathsf G\mathsf G^{\mathsf T} =
900 + 1600 = 2500$, sehingga $\mathbf m = (30, 40)\times0.50/2500 =
(0.006, 0.008)\,\unit{s.km^{-1}}$. (c) $\mathsf R = \frac{1}{2500}
\begin{pmatrix} 900 & 1200 \\ 1200 & 1600 \end{pmatrix} =
\begin{pmatrix} 0.36 & 0.48 \\ 0.48 & 0.64 \end{pmatrix}$, dengan jejak 1: satu
data menentukan satu kombinasi parameter. (d) $\mathsf R\,(0.01, 0) =
(0.0036, 0.0048)\,\unit{s.km^{-1}}$: anomali tersebar sepanjang sinar, lebih
banyak ke dalam sel tempat sinar lebih panjang.
\end{solution}

\begin{solution}{exo:b3:seismology-imaging:9}
(a) Secara linear: $\delta t = -(\delta v/v)L/v = -0.03\times200/9.0 =
\qty{-0.667}{s}$; secara eksak, $L/(1.03v) - L/v = 200/9.27 - 200/9.0 =
\qty{-0.647}{s}$. (b) Lintasan di dalam lempengan adalah $100/\cos\ang{40} =
\qty{130.5}{km}$: $-0.03\times130.5/9.0 = \qty{-0.435}{s}$ secara linear dan
$\qty{-0.422}{s}$ secara eksak. (c) Nilai eksaknya sama dengan nilai linear yang dibagi dengan
$1.03$: galat sebesar \qty{3}{\percent}, setara orde $\delta v/v$ itu sendiri, seperti
yang diharapkan dari teori orde pertama.
\end{solution}

\begin{solution}{exo:b3:seismology-imaging:10}
(a) Mode lebih tinggi ada selama bagian setengah-ruangnya meluruh, $c < \beta_2$; pada
batas potong $c = \beta_2$, sehingga $q_2 = 0$ dan $\tan(kHq_1) = 0$, dengan $kHq_1$ di
$[n\pi, n\pi + \pi/2)$: $kHq_1 = n\pi$. (b) Dengan $k = \omega/\beta_2$ dan
$q_1 = \sqrt{\beta_2^2/\beta_1^2 - 1} = \sqrt{\beta_2^2 - \beta_1^2}/\beta_1$,
$\omega_nH\sqrt{\beta_2^2 - \beta_1^2}/(\beta_1\beta_2) = n\pi$, yakni
rumusnya. (c) $\sqrt{4.5^2 - 3.5^2} = \sqrt8 = 2.828$; $\omega_1 =
\pi\times3.5\times4.5/(35\times2.828) = \qty{0.4998}{rad.s^{-1}}$, dengan periode
$2\pi/0.4998 = \qty{12.57}{s}$: pada periode yang lebih panjang hanya mode fundamental
yang ada. (d) Untuk $n = 0$ frekuensi batas potongnya nol: mode fundamental
ada pada semua periode. Pada periode pendek $kH \to \infty$ sedangkan $kHq_1$ tetap
di bawah $\pi/2$, sehingga $q_1 \to 0$ dan $c \to \beta_1$.
\end{solution}

\begin{solution}{exo:b3:seismology-imaging:11}
(a) $\mathsf G^{\mathsf T}(\mathsf G\mathsf G^{\mathsf T} + \varepsilon^2\mathsf I)
= (\mathsf G^{\mathsf T}\mathsf G + \varepsilon^2\mathsf I)\mathsf G^{\mathsf T}$;
kalikan dari kiri dengan $(\mathsf G^{\mathsf T}\mathsf G + \varepsilon^2
\mathsf I)^{-1}$ dan dari kanan dengan $(\mathsf G\mathsf G^{\mathsf T} +
\varepsilon^2\mathsf I)^{-1}$, keduanya invertibel untuk $\varepsilon > 0$.
(b) Di sini $\mathsf G\mathsf G^{\mathsf T} = \|\mathbf g\|^2 = 2500$, sebuah bilangan,
sehingga $\mathsf R_\varepsilon = \mathsf G^{\mathsf T}\mathsf G/(2500 + \varepsilon^2)$,
dengan jejak $2500/(2500 + \varepsilon^2)$. (c) Data prediksi adalah
$\mathbf g\cdot\mathbf m_\varepsilon = 0.50\times2500/(2500 + \varepsilon^2)$
dan residunya $0.50\,\varepsilon^2/(2500 + \varepsilon^2)$. Untuk
$\varepsilon^2 = 0$: jejak 1, residu 0; untuk 2500: jejak 0.5, residu
\qty{0.25}{s}; untuk \num{10000}: jejak 0.2, residu \qty{0.40}{s}. Redaman yang lebih besar
menghasilkan model yang lebih kecil, resolusi yang lebih buruk, dan ketakcocokan yang lebih besar.
\end{solution}

\begin{solution}{exo:b3:seismology-imaging:12}
(a) Dari latihan 5, $T = (2R/v)\sin(\Delta/2)$, maka $\mathrm dT/\mathrm d\Delta
= (R/v)\cos(\Delta/2)$. Tali busur yang membentuk sudut $\Delta$ di pusat adalah
alas segitiga sama kaki dengan sudut puncak $\Delta$, sehingga ia memotong jari-jari
di sumber pada $i_0 = (\pi - \Delta)/2$, dan $p = R\sin i_0/v =
(R/v)\cos(\Delta/2)$: sama. (b) \omterm{def:b3:seismology-imaging:ray-parameter}{Titik balik} terletak di tengah
tali busur, pada jarak $R\cos(\Delta_1/2)$ dari pusat; secara setara
$r_1 = p_1v$. (c) Di sini $p(\Delta)/p_1 = \cos(\Delta/2)/\cos(\Delta_1/2)$, dan
rumusnya memberikan $r_1 = R\exp\bigl(-\ln(1/\cos(\Delta_1/2))\bigr) =
R\cos(\Delta_1/2)$, seperti pada (b); kecepatannya adalah $r_1/p_1 = R\cos(\Delta_1/2)\big/
\bigl((R/v)\cos(\Delta_1/2)\bigr) = v$.
\end{solution}

\begin{solution}{pb:b3:seismology-imaging:1}
\textbf{1.} $i_0 = \arcsin 0.8 = \ang{53.13}$ dan $p = 0.8/5.0 =
\qty{0.16}{s.km^{-1}}$.

\textbf{2.} Sinar berbalik di tempat $v = 1/p = \qty{6.25}{km.s^{-1}}$, yakni pada
$z = (6.25 - 5.0)/0.05 = \qty{25}{km}$.

\textbf{3.} Dengan mendiferensialkan $\sin i = pv$ sepanjang sinar, $\cos i\,\mathrm di/
\mathrm ds = pg\,\mathrm dz/\mathrm ds = pg\cos i$, maka $\mathrm di/\mathrm ds =
pg$: arahnya berbelok dengan laju tetap, yakni kelengkungan lingkaran dengan jari-jari
$1/(pg) = 1/(0.16\times0.05) = \qty{125}{km}$. Di \omterm{def:b3:seismology-imaging:ray-parameter}{titik balik}, sinar
mendatar pada kedalaman \qty{25}{km} dan pusatnya terletak secara vertikal \qty{125}{km}
di atasnya. Pada $z = \qty{-100}{km}$, yakni $v_0/g$ di atas permukaan, tempat
kecepatan hasil ekstrapolasi akan nol.

\textbf{4.} Di sumber, jari-jari lingkaran tegak lurus terhadap sinar,
sehingga ia membentuk sudut $i_0$ dengan garis horizontal, dan pusatnya terletak
sejauh $125\cos i_0$ secara horizontal; berdasarkan simetri
$X = 2\times125\times\cos i_0 = 2\times125\times0.6 = \qty{150}{km}$, atau
$X = 2\sqrt{1 - p^2v_0^2}/(pg)$.

\textbf{5.} $\mathrm dT = \mathrm ds/v$ dengan $\mathrm ds = \mathrm di/(pg)$ dan
$v = \sin i/p$: $\mathrm dT = \mathrm di/(g\sin i)$. Selanjutnya
\[
  T = \frac2g\int_{i_0}^{\pi/2}\frac{\mathrm di}{\sin i} = \frac2g\ln\cot\frac{i_0}2 =
  \frac2g\operatorname{arccosh}\frac1{pv_0} .
\]
Di sini $\cot(\ang{26.57}) = 2$, maka
$T = 40\ln2 = \qty{27.73}{s}$.

\textbf{6.} Untuk $p = 0.15$: $pv_0 = 0.75$, $\sqrt{1 - 0.5625} = 0.6614$,
$X = 2\times0.6614/0.0075 = \qty{176.4}{km}$ dan $T = 40\operatorname{arccosh}
(1.3333) = 40\times0.7954 = \qty{31.81}{s}$. Selanjutnya $\Delta T/\Delta X =
4.08/26.4 = \qty{0.155}{s.km^{-1}}$, berada di antara 0.15 dan 0.16, sebagaimana
dituntut oleh $\mathrm dT/\mathrm dX = p$.

\textbf{7.} $\mu_1 = 2700\times3500^2 = \qty{33.1}{GPa}$ dan
$\mu_2 = 3300\times4500^2 = \qty{66.8}{GPa}$.

\textbf{8.} $k = 2\pi/(cT) = 2\pi/(3.99\times25) = \qty{0.06299}{rad.km^{-1}}$;
$q_1 = \sqrt{(3.99/3.5)^2 - 1} = 0.5474$; $q_2 = \sqrt{1 - (3.99/4.5)^2} = 0.4624$.

\textbf{9.} $\tan(kHq_1) = 66.825\times0.4624/(33.075\times0.5474) = 1.707$; mode fundamental
mengambil $kHq_1 = \arctan1.707 = 1.041$, sehingga
$H = 1.041/(0.06299\times0.5474) = \qty{30.2}{km}$.

\textbf{10.} Pada \qty{40}{s}: $k = 2\pi/170 = 0.03696$, $q_1 = 0.6888$,
$q_2 = 0.3287$, $\tan(kHq_1) = 0.964$, $kHq_1 = 0.7671$ dan
$H = 0.7671/(0.03696\times0.6888) = \qty{30.1}{km}$. Kedua hasil sesuai dalam batas
ketelitian kecepatan fase: kerak tebalnya kira-kira \qty{30}{km}.

\textbf{11.} $\omega = 0.2513$ dan $\qty{0.1571}{rad.s^{-1}}$, $k = 0.06299$ dan
$\qty{0.03696}{rad.km^{-1}}$: $U \approx 0.0942/0.02603 = \qty{3.62}{km.s^{-1}}$,
lebih kecil daripada kedua kecepatan fase, sebagaimana mestinya bila $c$ meningkat seiring
periode.

\textbf{12.} $\lambda = 3.99\times25 = \qty{99.8}{km}$ dan $4.25\times40 =
\qty{170}{km}$. Mode fundamental mengambil sampel hingga sebagian kecil dari
panjang gelombangnya: gelombang 25-s sebagian besar terletak di kerak, gelombang 40-s mencapai jauh
ke dalam mantel, sehingga kecepatannya lebih dekat ke $\beta_2$. Dua periode yang menghasilkan
$H$ yang sama menjadi pemeriksaan bahwa model dua lapisan memadai.

\textbf{13.} Dengan kolom untuk sel 1 sampai 4 dan baris A sampai F, dalam
satuan \qty{1000}{km} ($\sqrt2\times100/1000 = 0.1414$):
\[
  \mathsf G = \begin{pmatrix}
  0.1 & 0.1 & 0 & 0 \\ 0 & 0 & 0.1 & 0.1 \\ 0.1 & 0 & 0.1 & 0 \\
  0 & 0.1 & 0 & 0.1 \\ 0.1414 & 0 & 0 & 0.1414 \\ 0 & 0.1414 & 0.1414 & 0
  \end{pmatrix}.
\]

\textbf{14.} $\mathsf C_d^{-1} = \operatorname{diag}(400, 400, 400, 400, 100,
100)$. Sinar baris atau kolom menyumbang $0.1^2\times400 = 4$, sinar diagonal
$0.1414^2\times100 = 2$. Sel 1 dilintasi oleh A, C (masing-masing 4) dan E (2): 10. Sel
1 dan 2 berbagi sinar A: 4; sel 1 dan 3 berbagi C: 4; sel 1 dan 4 berbagi E: 2; dan
seterusnya berdasarkan simetri, yang menghasilkan keempat baris yang dinyatakan.

\textbf{15.} Sinar baris atau kolom menyumbang $0.1\times400\,\delta t = 40\,
\delta t$ ke setiap selnya, sinar diagonal $0.1414\times100\,\delta t =
14.14\,\delta t$. Sel 1: $40(0.16 - 0.36) + 14.14\times(-0.15) = -10.12$; sel
2: $40(0.16 + 0.11) + 14.14\times(-0.27) = 6.98$; sel 3: $40(-0.48 - 0.36) +
14.14\times(-0.27) = -37.42$; sel 4: $40(-0.48 + 0.11) + 14.14\times(-0.15) =
-16.92$.

\textbf{16.} Untuk $(1,1,1,1)$, setiap baris $\mathsf N$ berjumlah 20. Untuk
$(1,-1,-1,1)$, baris 1 memberikan $10 - 4 - 4 + 2 = 4$ dan baris lainnya mengikuti
pola yang sama. Untuk $(1,1,-1,-1)$, $10 + 4 - 4 - 2 = 8$, dan untuk $(1,-1,1,-1)$,
$10 - 4 + 4 - 2 = 8$. Vektor-vektor tersebut ortogonal dengan norma 2, sehingga
$\mathsf N^{-1} = \sum_k \mathbf u_k\mathbf u_k^{\mathsf T}/(4\lambda_k)$; entri
$(1,1)$-nya adalah $\tfrac14(\tfrac1{20} + \tfrac14 + \tfrac18 + \tfrac18) =
\tfrac{11}{80}$, entri $(1,2)$-nya $\tfrac14(\tfrac1{20} - \tfrac14 + \tfrac18
- \tfrac18) = -\tfrac4{80}$, entri $(1,4)$-nya $\tfrac14(\tfrac1{20} +
\tfrac14 - \tfrac18 - \tfrac18) = \tfrac1{80}$, dan seterusnya.

\textbf{17.} $\hat{\mathbf m} = \mathsf N^{-1}\mathbf b$:
$\hat m_1 = (11\times(-10.121) - 4\times6.982 - 4\times(-37.418) - 16.921)/80
= -0.081$; dengan cara serupa $\hat m_2 = 1.844$, $\hat m_3 = -3.706$ dan
$\hat m_4 = \qty{-0.931}{ms.km^{-1}}$.

\textbf{18.} $\delta v/v_0 = -v_0\,\delta s = -8.0\times10^{-3}\,\hat m$, dalam persen
$-0.8\,\hat m$: sel 1 $+0.07$, sel 2 $-1.48$, sel 3 $+2.96$, sel 4
$\qty{+0.75}{\percent}$.

\textbf{19.} Setiap entri diagonal $\mathsf N^{-1}$ bernilai $11/80$: setiap sel
memiliki galat baku $\sqrt{0.1375} = \qty{0.371}{ms.km^{-1}}$, yaitu
$\qty{0.30}{\percent}$. Korelasi sel 1 dan 2 adalah $-4/11 = -0.36$:
keduanya berbagi sinar A, dan perlambatan yang diberikan pada satu sel dapat
dipertukarkan dengan sel lainnya. Sel 1 dan 4, yang hanya berbagi diagonal E
yang kurang presisi, berkorelasi lemah dan positif, $+1/11 = +0.09$.

\textbf{20.} Residu prediksi adalah $0.176$, $-0.464$, $-0.379$, $0.091$,
$-0.143$ dan $\qty{-0.263}{s}$; residu setelah pencocokan adalah $-0.016$,
$-0.016$, $0.019$, $0.019$, $-0.007$ dan $\qty{-0.007}{s}$, dan
$\chi^2 = 0.50$ untuk $6 - 4 = 2$ derajat kebebasan: konsisten dengan galat
yang dinyatakan, bila ada, sedikit di sisi kecil.

\textbf{21.} Vektor eigen adalah milik $\mathsf N$ dan nilai eigennya
$\lambda/(\lambda + 4)$: $20/24 = 5/6$, $4/8 = 1/2$, $8/12 = 2/3$ dua kali. Seperti
pada soal 16, $R_{11} = \tfrac14(\tfrac56 + \tfrac12 + \tfrac23 + \tfrac23) =
\tfrac23$; $R_{12} = \tfrac14(\tfrac56 - \tfrac12 + \tfrac23 - \tfrac23) =
\tfrac1{12}$; $R_{13} = \tfrac14(\tfrac56 - \tfrac12 - \tfrac23 + \tfrac23) =
\tfrac1{12}$; $R_{14} = \tfrac14(\tfrac56 + \tfrac12 - \tfrac23 - \tfrac23) = 0$.

\textbf{22.} Proyeksikan $\mathbf b$ pada vektor-vektor eigen: $\mathbf b\cdot\mathbf u_k/4
= -14.37$, $0.849$, $12.80$ dan $-9.40$ untuk nilai eigen 20, 4, 8 dan 8.
Taksiran teredam membaginya dengan $\lambda + 4$ alih-alih $\lambda$; untuk
sel 3 (komponen $1, -1, -1, 1$),
$m_3 = -14.37/24 - 0.849/8 - 12.80/12 - 9.40/12 = \qty{-2.555}{ms.km^{-1}}$,
anomali kecepatan sebesar $\qty{+2.04}{\percent}$. Peredaman menarik setiap
komponen mendekati nol, paling kuat pada komponen yang paling lemah terikat.

\textbf{23.} Model tanpa redaman, yang mencocokkan data sampai batas galatnya
($\chi^2 = 0.50$ untuk 2 derajat kebebasan), sedangkan model teredam
mencocokkan data dengan sangat buruk ($15.8$). Peredaman hanya diperlukan bila
$\mathsf N$ singular atau mendekati singular; di sini nilai eigen terkecilnya, 4,
tidak kecil, taksirannya stabil dan galatnya moderat.

\textbf{24.} Baris 3 adalah $(1/12, 0, 2/3, 1/12)$: taksiran teredam untuk sel 3
adalah dua pertiga nilai sejatinya ditambah seperduabelas dari sel 1 dan 4,
tetangganya sepanjang sinar C dan sinar B, dan tidak ada apa pun dari sel 2.

\textbf{25.} Sel 3, sel bawah di sebelah kiri, cepat sekitar
\qty{3}{\percent}; sel 2 lambat. Lempengan litosfer samudra yang dingin dan padat
yang tenggelam di bawah benua lebih cepat daripada mantel di sekitarnya, sehingga
sel cepat, dalam dan di sisi samudra dari K, adalah lempengan tersebut; sel lambat
di atasnya di sisi lain terletak di tempat baji mantel panas.

\textbf{26.} Dari soal 18 dan 19, anomali kecepatan gelombang P pada sel lempengan adalah
\textbf{\qty[text-series-to-math=true, uncertainty-mode=separate]{3.0 \pm 0.3}{\percent}}.
\end{solution}
